\section{Earlier controls and external positivity proposals}
\label{sec:controls}

These results helped formulate the physical question, but each has its own
action and domain. They are not alternative derivations of the curved
energy theorem. In particular an instability proved for a flat surrogate
is not a growing-mode theorem for the curved background.

\subsection{Positive spatial pairings do not ensure stable wave energy}
For two complex scalar components on $\RR$, fix $b>0$ and a smooth real
static function $r(x)$. Define
\begin{equation}
 K=\begin{pmatrix}r&1\\1&0\end{pmatrix},\quad
 B=\begin{pmatrix}0&0\\r'&0\end{pmatrix},\quad
 U=\begin{pmatrix}0&-b\\0&br\end{pmatrix},\quad
 D=-I\partial_x^2-B\partial_x+U.
 \label{eq:twowave}
\end{equation}
The declared wave equation is $u_{tt}+Du=0$, from the action
$\frac12\int(u_t^TKu_t-u_x^TKu_x+u^T\diag(0,b)u)$.
This is the observed two-field model, without an I1B identification
\cite{StaticPairing}.

Formal symmetry in a Hermitian multiplication pairing $W$ requires
$2W'=B^TW+WB$ and $WU-U^TW+W''-(B^TW)'=0$.
Solving both equations, including an imaginary off-diagonal entry before
it is forced to vanish, gives every positive solution:
\begin{equation}
 W=s\begin{pmatrix}
 c^2+\beta r+r^2/4&\beta+r/2\\\beta+r/2&1
 \end{pmatrix},\quad
 s>0,\ c>0,\ |\beta|<c,\quad
 r''=\frac b2(r^2-4c^2).\label{eq:pairingclass}
\end{equation}
For $s=1$, $\beta=0$, the factor
$F=\bigl(\begin{smallmatrix}c&0\\r/2&1\end{smallmatrix}\bigr)$
satisfies $W=F^TF$. Conjugating and applying a constant orthogonal
rotation gives scalar channels
\begin{equation}
 D_-=-\partial_x^2+b(r/2-c),\qquad
 D_+=-\partial_x^2+b(r/2+c).\label{eq:scalar-channels}
\end{equation}
Every other allowed pairing becomes a positive constant weight on these
same channels. For bounded smooth $r$, the profile equation and its first
integral bound the derivatives needed for this similarity, giving a
closed realization from the scalar $H^2(\RR)$ domains.

\begin{proposition}[Bounded-profile energy criterion]
Within \eqref{eq:pairingclass} on the whole line, $D\ge0$ precisely when
$r\equiv2c$.
\end{proposition}
\begin{proof}
A bounded profile cannot exceed $2c$: on any component where it does,
$r''>0$, and convexity contradicts either the endpoint values or
boundedness on a half-line or the whole line. Thus $b(r/2-c)\le0$.
Unless $r=2c$ identically, it is strictly negative on some interval with
integral $-\delta<0$. A trial function equal to one there and tapering
to zero over lengths $L$ has kinetic cost $2/L$ and potential at most
$-\delta$. Large $L$ and smooth approximation give a negative form
value. At $r=2c$ both scalar potentials are nonnegative.
\end{proof}

For $k^2=bc/2$, the exact nonconstant profile
$r=2c-6c\sech^2(kx)$ still has uniformly positive $W$.
Its $D_-$ eigenfunctions $\sech^2(kx)$ and
$\sinh(kx)\sech^2(kx)$ have eigenvalues $-4k^2$ and $-k^2$.
They give growing wave solutions with rates $2k$ and $k$ on the stated
domain. A positive conserved phase-space pairing cannot contain a
nonzero generator eigenvector with positive real growth rate: if
$Av=\gamma v$ and $A^\dagger G+GA=0$, then
$2\gamma v^\dagger Gv=0$. This is a wave statement, distinct from
positivity of the spatial weight and from the original indefinite
canonical action.

\subsection{Flat I1B surrogates and their correction}
\label{sec:flat}
The flat constant-coefficient completion used the same selected Shiab
but set the ambient connection and curvature to zero. On fixed flat
geometry, the radial restriction $T=\phi(t)\Phi_1$ omits normal bivector
equations requiring $\phi'=0$. Adding one isotropic bivector amplitude
$v$ still gives necessary equations
\begin{align}
 \kappa\phi+312\phi^2-\tfrac{260}{3}v^2&=0,\nonumber\\
 v'&=\tfrac4{33}v^2,\qquad
 -11\phi'-(\kappa+\tfrac{568}{3}\phi)v=0.\label{eq:flatclosure}
\end{align}
On an interval with $v\ne0$, differentiating the first equation forces
$118976\phi^2+816\kappa\phi+\kappa^2=0$.
Continuity then makes $\phi$ constant, contradicting the $v'$ equation.
Thus this small nonlinear restriction supplies no varying solution.

A larger linearized metric/distortion sector does have a growing mode
in that surrogate. With diagonal spatial metric amplitude $\psi$, let
$(a,b,d)$ be time, horizontal-spatial, and vertical grade-one coefficients,
and $(u,w)$ their two bivector companions. The metric condition is
$(b+5d)''=0$, and the coupled characteristic determinant is
\begin{equation}
 21600\kappa^2z^4(113z^2-17\kappa^2).
 \label{eq:flatdet}
\end{equation}
For $\kappa\ne0$ an exact mode has
\begin{equation}
 d=e^{\gamma t},\quad b=-5d,\quad a=0,\quad
 u=-5d'/\kappa,\quad w=7d'/\kappa,\quad
 \psi=135d/(17\kappa),\quad
 \gamma=|\kappa|\sqrt{17/113}.\label{eq:flatmode}
\end{equation}
The certificate checks all ten metric rows and all 28 distortion rows
in the closed label block; trace selection excludes omitted receivers.
On the zero-affine branch its reduced energy is
$565(d')^2/(2\kappa)-85\kappa d^2/2$.
The separate distortion-only exponent $\kappa/11$ is removed by metric
coupling. Neither calculation is transferred to the curved equilibrium.

The associated co-moving flat constant-background experiment also has
an exact exclusion. Its metric condition $L=0$, together with distortion
scaling $3Q+\kappa N=0$, forces $Q=N=0$ at nonzero $\kappa$.
On diagonal grade-one fields $N=\sum t_\mu^2$.
Adding the two bivector companions makes $N$ indefinite, but the exact
necessary stationary ideal, rescaled to unit coupling, is
\begin{equation}
 \langle\partial_a(Q+N/2),\partial_b(Q+N/2),
 \partial_d(Q+N/2),\partial_u(Q+N/2),\partial_w(Q+N/2),N\rangle
 =\langle a,b,d,u,w\rangle.\label{eq:flatideal}
\end{equation}
This rejects that five-field nonzero constant family \cite{FlatConstant}.
Its $-8L$ metric variation belongs to the declared co-moving completion;
\cref{sec:background} explains why it is not the independent $g$
variation of the canonical metric bundle. The constant-background
exclusion and growing mode remain exact results on their named object,
with no canonical-curved or physical GU conclusion \cite{FlatScalar}.

\subsection{The scope of the original objection}
Nguyen--Polya's response is dated February 23, 2021 and discusses the
2013 lecture with 2020 follow-up material as a pre-quantum classical
proposal \cite[pp.~1--2]{NP}. Weinstein's written draft is dated
April 1, 2021 \cite{GU}; later formulas cannot be treated as text the
February authors overlooked. The response's Section 3.1 combines a
real-structure objection with an inference from complexification to
nonunitarity or two-sided unbounded energy.
An abstract real vector-space isomorphism is different from an
identification preserving representation, real structure, and operations.
Even if a real identification needs correction, complexifying notation
does not itself declare all complex components independent real physical
fields. The source-phased slice in \eqref{eq:packets} is one reason those
steps must be distinguished.

A charitable dynamical reading asks about a propagating Yang--Mills
sector with an indefinite invariant kinetic form. That is a substantive
conditional concern, but its action and physical modes need an
applicability proof. The following comparator makes those assumptions
explicit.

\section{The Yang--Mills obstruction with its hypotheses exposed}
\label{sec:ym}

Let $G$ be a connected Lie group whose real Lie algebra $\mathfrak g$
has a nondegenerate, $\operatorname{Ad}(G)$-invariant symmetric bilinear
form $\beta$. On
$\mathbb R_t\times\mathbb T^d$, $d\geq2$, use the flat metric of signature
$(-,+,\ldots,+)$, a trivial bundle, real connection components and
\begin{equation}
 S_{\YM}=-\frac1{4g^2}\int\beta(F_{\mu\nu},F^{\mu\nu})\,dt\,dx,
 \qquad g^2>0. \label{eq:ym}
\end{equation}
With $E_i=F_{0i}$, the energy and Gauss constraint are
\begin{align}
 H&=\frac1{2g^2}\int_{\mathbb T^d}
 \left(\sum_i\beta(E_i,E_i)+
 \sum_{i<j}\beta(F_{ij},F_{ij})\right)dx, \label{eq:H}\\
 D_iE_i&=0. \label{eq:Gauss}
\end{align}
Nondegeneracy of $\beta$ allows the usual identification of canonical
momentum with $g^{-2}\beta(E,\cdot)$. All these assumptions belong to this
comparator, not to GU by default.

\begin{proposition}[Negative energy survives the ordinary Gauss constraint]
\label{prop:negative}
If some $X\in\mathfrak g$ has $\beta(X,X)<0$, the Hamiltonian
\eqref{eq:H} is unbounded below on smooth Gauss-compatible initial data,
even after ordinary gauge equivalence. If $\beta$ has both signs, the
Hamiltonian is also unbounded above.
\end{proposition}
\begin{proof}
Choose a nonzero smooth divergence-free real vector field $u$ on the
torus and take
\[
 A_i=0,\qquad E_i=\lambda X u_i.
\]
Then $F_{ij}=0$ and $D_iE_i=\lambda X\partial_i u_i=0$, whereas
\[
 H=\frac{\lambda^2\beta(X,X)}{2g^2}
       \int_{\mathbb T^d}|u|^2dx\longrightarrow-\infty.
\]
Gauge transformations preserve $H$ by invariance of $\beta$, so taking
the gauge quotient cannot remove these values. Using a positive direction
instead proves the assertion about the upper bound. For a nonconstant
transverse example one may choose $u_1=\sin x^2$ and $u_i=0$ for $i\ne1$.
\end{proof}

This is an exact classical result, not a finite-difference instability
or a proof about an unspecified quantum domain. It demonstrates why the
negative internal direction is not merely the removable timelike
polarization of ordinary compact Yang--Mills theory. On a one-color
abelian subsector, $[X,X]=0$, so the negative sign is already visible in
an ordinary transverse Maxwell mode.

For example, if a complex Lie algebra carries a nondegenerate complex
symmetric invariant form $b$, its underlying real vector space admits
\[
 \beta_c(X,Y)=\operatorname{Re}\bigl(c\,b(X,Y)\bigr),\qquad c\ne0.
\]
It is real nondegenerate: vanishing against both $Y$ and $iY$ forces
$b(X,Y)=0$ for every $Y$. Moreover,
$\beta_c(iX,iX)=-\beta_c(X,X)$, so its signature is split. This explains
the kinetic-sign problem when this form and these real independent fields
are actually used in \eqref{eq:ym}. It says nothing about a theory with
a different action or physical reduction.

The usual quantum alternative is equally concrete at the linear level.
A negative transverse oscillator quantized on positive $L^2(\mathbb R)$
has
\[
 \widehat H=-\tfrac12(\widehat p^{\,2}+\omega^2\widehat q^{\,2}),
 \qquad \operatorname{spec}(\widehat H)
 =\{-\omega(n+\tfrac12):n\geq0\}
\]
in units $\hbar=1$. It is self-adjoint and generates unitary time
evolution, but is not bounded below. Exchanging the roles of annihilation
and creation operators to reverse the ladder gives a negative one-particle
norm if the same adjoint and canonical commutator are retained. This
explains the conventional dilemma; it is not a classification of all
possible contours, representations or constrained quantum theories.
Witten's introductory comparison makes the corresponding action-dependent
point \cite[report pp.~2--3]{WittenReport}.

\section{What Witten's construction does and does not say}
\label{sec:witten}

Witten studies three-dimensional Chern--Simons theory for the
complexification $G_{\mathbb C}$ of a compact group $G$, not a propagating
Yang--Mills curvature-square theory. His action pairs a
complex connection $\mathcal A$ with its complex conjugate in two
Chern--Simons terms, with couplings $t=k+is$ and
$\widetilde t=k-is$. In the real-$s$ branch relevant here, $k$ is
quantized and $s$ is real, with $t,\widetilde t\ne0$ for the displayed
connection formulas (for example $k=s=1$); the physical reality condition and the
choice of polarization are part of the construction. No assertion about
arbitrary complex couplings is needed. The alternative imaginary-$s$
positive construction is explicit only in genus one with restricted
parameters; no all-genus positivity assertion from that branch is used.

In canonical quantization on $\Sigma\times\mathbb R$, with $\Sigma$
a closed oriented surface, the time component of the connection imposes
flatness. The phase space is a moduli space of flat complex connections,
not the unconstrained space of transverse Yang--Mills waves. The
Hamiltonian of this pure theory is a constraint and vanishes on reduced
physical states. Thus the time evolution of this fixed, source-free
closed-surface problem is the identity. A zero Hamiltonian is bounded
below, but this does not imply that the state space is empty or that all
topological observables are trivial. Boundary terms, puncture/source
insertions and chosen boundary dynamics require a separate analysis;
one may not export $H=0$ to those problems or to a propagating theory.

The journal identity is \cite{Witten1991}; the precise theorem used here concerns a \emph{connection on a Hilbert
bundle over Teichm\"uller space}, the parameter space of auxiliary complex
structures on $\Sigma$. The fiber is described as $L^2(\mathcal M,L^k)$,
where $\mathcal M$ is the compact-group moduli space used in the chosen
polarization and $L^k$ its quantization line bundle. For real $s$, the
specified corrected connection $\delta^Q$ is unitary and projectively
flat, with central curvature of type $(1,1).$ The pairing has the form
$\int_{\mathcal M} H(\chi,\psi)\,d\mu$, where $H$ incorporates the
regularized Laplacian determinant; conjugation by $H^{1/2}$ identifies it
with the ordinary positive $L^2$ pairing. This is unitary
\emph{parameter-space transport}, not a theorem identifying a
four-dimensional physical-time Hamiltonian. Projective flatness also
retains a scalar-phase issue; it is not literal flatness.

There is a further rigor qualification. Witten gives heuristic reasoning
for the quantum correction to the connection, then proves the stated
properties of the selected formulas. Moduli-space singularities,
noncompactness and the definition of general functional integrals are
not thereby eliminated. In particular, the identification of the reduced
state space uses the dense stable locus and the corresponding
Narasimhan--Seshadri description; it is not an unrestricted assertion
about every singular stratum. The paper supplies a mathematical construction
and a substantive counterexample to the group-only slogan in its stated
quantization setting; it is not a complete axiomatic construction of
interacting QFT for every complex group. The positive pairing and the
closed-surface zero-Hamiltonian argument are separate ingredients.
Neither should be replaced by the bare word ``unitary.''

All detailed formula locators in this paragraph refer to the
44-page IASSNS-HEP-89/65 report: \S3, p.~8 (closed surface);
\S4, pp.~19--21 (correction and theorem); \S4.1, pp.~22--23,
Eqs.~(4.12) and (4.16) (pairings) \cite{WittenReport}.
The journal bibliographic identity is verified, but its typesetting was
not independently inspected; no journal page is used as a formula locator.


\subsection{Perturbative noncompact Chern--Simons and the GU transfer}
Bar-Natan--Witten study perturbation theory with a noncompact gauge group
and a gauge-fixed kinetic construction using an auxiliary positive form
\cite{BNW}. Its gauge fixing and operator are part of that construction;
changing the pairing does not leave an arbitrary action and domain
unchanged. Neither this perturbative result nor Witten's
parameter-space theorem identifies an interacting four-dimensional GU
Hamiltonian.

In the source, $I_1^B$, its printed residual
$\Upsilon^B=\cS(F_A)+*\kappa T$, and $I_2^B=\|\Upsilon^B\|^2$ are
separate variational owners. For a residual-square functional
$S(\phi)=Q_\phi(R(\phi),R(\phi))$ at $R(\phi_0)=0$,
\[
 D^2S_{\phi_0}(v,w)=2Q_{\phi_0}(DR_{\phi_0}v,DR_{\phi_0}w).
\]
An indefinite ambient form alone does not determine its restriction to
the derivative image. Nor does the symbol $\|\cdot\|^2$ prove positivity.
The direct I1B energy calculation in this paper differentiates its own
action and solves its own primary equations; it requires neither
comparator transfer.

Accordingly the group-only inference does not establish the GU-specific
energy conclusion. The ordinary Yang--Mills proposition establishes its
conditional obstruction on \eqref{eq:ym}. The I1B result establishes a
separate obstruction on \eqref{eq:action}. Citing complex Chern--Simons
alone supplies no positive physical completion of either propagating
model. These judgments bind the stated objects, not a critic or a theory
as an undifferentiated whole.

\subsection{Scalar and composite-observable donors}
The dated source intake also inspected four-derivative scalar and
composite-observable proposals \cite{ScalarIntake}. These are independent
models. In the fixed scalar basis $Q=\Box\sigma$, $X=(\partial\sigma)^2$,
a complete-square diagnostic for $aQ^2+bQX+cX^2$, $a\ne0$, is
$4ac-b^2=0$. In the normalization
$-Q^2/2+\lambda_3XQ+\lambda_4X^2$, it becomes
$\lambda_4=-\lambda_3^2/2$, with RG tangency
$\beta_4+\lambda_3\beta_3=0$ on the branch.
Passing that test does not prove closure of a scalar embedding under
normal Euler equations.

Anderson--Bateman--Herzog--Turok's inspected version gives three-loop
renormalization identities and a nonexceptional off-shell infrared
argument \cite{ABHT}; its on-shell inclusive probability problem is
separate. The cubic momentum polynomial reduces under conservation to
$2(p\cdot q)^2-2p^2q^2$, providing the soft-momentum factors needed in
that model. A transfer must retain that interaction grammar, not only
a conformal resemblance.
There is also a bounded displayed-coefficient correction: from
\[
 L=-\Upsilon\Box\Omega-\frac g6\Omega^2\Upsilon^2
\]
one obtains $\Upsilon=-3\Box\Omega/(g\Omega^2)$ and
$L_{\mathrm{eff}}=3(\Box\Omega/\Omega)^2/(2g)$ for
$g\Omega\ne0$. The inspected Eq.~(5.4) instead prints $1/(2g)$.
The source intake checks the PDF, HTML, and original TeX identity.
This discrepancy does not negate the separate loop calculation.
Bateman--Turok's tree-level ghost-parity result \cite{BT} does not
by itself establish loop positivity for a GU sector.

Amaral--Lemes' inspected second version constructs a cocycle in an
enlarged background-equivariant theory and a positive leading one-loop
composite spectral density \cite{AL}. BRST closure and positivity of the
full interacting operator are different conclusions.
Asorey and collaborators analyze a subtracted composite spectral
representation and a bubble-resummed bound-state correlator of
conjugate-mass fields \cite{AKPS}. Continuum density, isolated pole
residue, and contact subtractions must be checked separately; this is
not positivity of every Schwinger function. These exact-version source
records provide possible methods and tests, without a native GU transfer
or a claim about newer literature.

\section{Separate cutoff quantum and gauge--matter controls}
\label{sec:quantum}

The merged research corpus includes a distinct variational and PDE
sequence. This section states its current yield with its own objects;
none is used as a premise of \cref{thm:main}.
Take the normalized three-torus, cube cutoff
$K_N=\{k\in\mathbb Z^3:\|k\|_\infty\le N\}$, fixed $m,g>0$, and
$\omega_k=(m^2+|k|^2)^{1/2}$.
The positive Wick-square control has free Gaussian form $q_0$ and
\begin{equation}
 C_N=\sum_{k\in K_N}(2\omega_k)^{-1},\quad
 W_N=\int_{\TT^3}(\phi_N^2-3C_N)^2,\quad
 E_N=\inf_{\|\psi\|=1}\{q_0[\psi]+g\mathbb E_{|\psi|^2}W_N\}.
 \label{eq:quantum-control}
\end{equation}
This positive shifted polynomial convention must not be confused with
an unshifted normal-ordered quartic Hamiltonian.

\subsection{From sign texture to the quartic exponent}
For a real degree-$N$ field $f$, set
$\delta=\int(f^2-1)^2$ and let a fixed-ratio Fourier shell of
$\operatorname{sgn}f$ carry $L^2$ mass at least $\eta>0$.
Since $\|f-\operatorname{sgn}f\|_2^2\le\delta$, a small defect forces
$\|\nabla f\|_2^2\ge\alpha^2\eta N^2/4$.
On $B=\{|f|<1/2\}$, its measure is at most $16\delta/9$ and
Bernstein--H\"older gives gradient cost at most $CN^2\sqrt\delta$.
On $B^c$, differentiating $f^2-1$ bounds it by $4N^2\delta$.
After canceling $N^2$,
\[
 \alpha^2\eta/4\le C\sqrt\delta+4\delta.
\]
Thus the defect has a strictly positive cutoff-independent lower bound
under the fixed-shell hypothesis. A high-frequency sine has defect
$3/8$ and sign-shell mass $8/\pi^2$, so the exponent zero is sharp.
The probability-level lifting of this estimate, combined with the
low-shell Fisher bound, proves
\begin{equation}
 E_N=\Theta_g(N^4).\label{eq:N4}
\end{equation}
The vacuum supplies the matching upper bound $6gC_N^2$.
The earlier lamellar exclusions and $N^{-4/3}$ defect estimate remain
valid but are superseded, on the same fixed-shell problem, by this
stronger result \cite{QuarticControl}.
Equation \eqref{eq:N4} identifies an exponent, not the unrestricted
coefficient or a bounded-error recentering.

\subsection{A class coefficient and an unrestricted descent}
The cube geometry matters:
\[
 C_N/N^2\longrightarrow c_C=12\log\frac{1+\sqrt3}{\sqrt2}-\pi.
\]
For a continuum squeeze profile $s$ on $Q=[-1,1]^3$, define
\[
 A[s]=\frac12\int_Q\frac{s(x)}{|x|}\dd x,\qquad
 F[s]=\frac14\int_Q|x|(s+s^{-1}-2)\dd x.
\]
At fixed variance the unique free-cost minimizer is
$s_\kappa(x)=|x|/\sqrt{|x|^2+\kappa}$.
After optimizing the constant mean, the stationary diagonal Gaussian
coefficient is the minimum of $F_*(A)+6gA(2c_C-A)$.
Its positive parameter solves
$\kappa_g=24g(c_C-A(\kappa_g))$.
This improves the uniform-squeeze class for every $g>0$.
Finite-cutoff profiling supplies a global minimum within that Gaussian
class, denoted $\lambda_N^{\mathrm{prof}}$ \cite{ProfileControl}.
Moment-matched Gaussian comparison extends the same minimum to the
translation-invariant finite-Fisher class with nonnegative integrated
Wick defect $\int(4h\mu_3+\mu_4-3v^2)\ge0$
\cite{NonGaussianControl}.

The sign condition is restrictive. At the Gaussian minimum, the residual
has precisely third and fourth Gaussian chaos:
\begin{align}
 (H_N-\lambda_N^{\mathrm{prof}})G_N&=R_NG_N,\nonumber\\
 R_N&=g\int\left[:\eta^4:_{\Gamma_N}
                      +4h_N:\eta^3:_{\Gamma_N}\right],\nonumber\\
 \sigma_N^2&=g^2\left(96h_N^2\int\Gamma_N^3
                              +24\int\Gamma_N^4\right).
 \label{eq:ritz}
\end{align}
Positive Fourier convolution counts give
$\sigma_N=\Theta_g(N^{5/2})$.
The two-dimensional Ritz matrix on $G_N$ and $R_NG_N/\sigma_N$ has
energy drop
\[
 \Delta_N=\frac{\sqrt{\delta_N^2+4\sigma_N^2}-\delta_N}{2},
 \qquad |\delta_N|\le C_g\sigma_N,
 \qquad \Delta_N=\Theta_g(N^{5/2}).
\]
Thus the profiled Gaussian misses the unrestricted translation-invariant
bottom by at least this divergent amount. Its value plus an
$o(N^{5/2})$ shift cannot recenter that bottom to bounded error.
The drop is still $o(N^4)$ and does not decide the unrestricted leading
coefficient \cite{RitzControl}.

The latest result covers every finite-second-moment location law $\pi$
on the zero-mode coordinate with the same profiled covariance $S_N$.
Fisher convexity and the moment lower bound have exact gap
\[
 U-L=\frac{\omega_0\operatorname{Var}_\pi A}
 {s_{N,0}(s_{N,0}+\operatorname{Var}_\pi A)}
 \le\frac{\omega_0}{s_{N,0}}=O_g(N).
\]
Linearity of potential expectation then yields
\[
 \lambda_N^{\mathrm{prof}}-O_g(N)
 \le\inf_\pi Q_N(\pi)\le\lambda_N^{\mathrm{prof}}.
\]
Arbitrarily many zero-mode locations with this common covariance retain
the same leading class coefficient. Covariance heterogeneity, spatial
textures, and the Ritz direction above are outside the theorem
\cite{MixtureControl}.

\subsection{Exact cancellation without a global dispersive theorem}
The separate temporal-gauge Maxwell--matter line studies opposite charges
$\pm e$ on a torus. Let $a(t)$ be harmonic holonomy,
$M=\sum_k(|\phi_{+,k}|^2+|\phi_{-,k}|^2)$, and
$D=\sum_k k(|\phi_{+,k}|^2-|\phi_{-,k}|^2)$.
The exact bare Fourier energy satisfies
\begin{equation}
 E_0=H_a+ea\cdot D-\tfrac12e^2|a|^2M,\qquad
 E_0'=ea\cdot D'-\tfrac12e^2|a|^2M'.\label{eq:bareenergy}
\end{equation}
It removes $\dot a$ from this linear harmonic calculation, but for $m>0$
the bound still spends $2|e||a|+(e^2/m)|a|^2$.
Earlier current and radial normal forms, covariant integration by parts,
and moving analytic radii must retain their remaining budgets.
Static flat connections show that a raw potential-size radius cost is
not always dynamically necessary. Time-dependent periodic controls show
that positive conserved energy and Gauss neutrality alone need not make
harmonic electric total variation finite.

For a nonzero lattice vector $k$ perpendicular to the holonomy plane,
\begin{equation}
 a(t)=A(\cos\Omega t,\sin\Omega t,0),\qquad
 \phi_+=\phi_-=R e^{i\omega t}e^{ik\cdot x},\label{eq:periodic}
\end{equation}
with $\Omega=\sqrt2eR$ and
$\omega^2=m^2+|k|^2+e^2A^2$, solves the stated coupled control.
Opposite Gauss densities cancel; the currents add to $-2e^2R^2a$.
It has no homogeneous matter mode, but $|\dot a|$ is nonzero constant.
Therefore excluding only $k=0$ does not produce dispersion. If the two
frequencies are commensurate the full fields are periodic; in general
the holonomy is periodic and the fields are bounded quasiperiodic.
The energy and nonintegrable-variation conclusions hold in either case.
This does not reject scattering under stronger localization or
nonresonance hypotheses \cite{MixtureControl}.

None of \eqref{eq:N4}--\eqref{eq:periodic} supplies a source-owned
continuum Hamiltonian, interacting BRST closure, physical quotient,
or a global evolution theorem for the curved I1B action.
